\documentclass[11pt,a4paper]{article}
\usepackage[utf8]{inputenc}
\usepackage[T1]{fontenc}
\usepackage{lmodern}
\usepackage[margin=2cm,top=2.2cm,bottom=2.2cm]{geometry}
\usepackage{amsmath,amssymb}
\usepackage{booktabs}
\usepackage{tabularx}
\usepackage{enumitem}
\usepackage{fancyhdr}
\usepackage{titlesec}
\usepackage{xcolor}
\usepackage{listings}
\usepackage{microtype}
\usepackage[colorlinks=true,linkcolor=blue!70!black,citecolor=blue!70!black,urlcolor=blue!70!black]{hyperref}

% Header and Footer styling
\pagestyle{fancy}
\fancyhf{}
\fancyhead[L]{\small\textsf{\textbf{HashiCorp Vault Engineering Guide} \textbar\ Architecture, Operations \& HVAC}}
\fancyhead[R]{\small\textsf{Page \thepage}}
\fancyfoot[L]{\small\textsf{Guardian Layer Security Specification}}
\fancyfoot[R]{\small\textsf{NIST SP 800-57 / Raft HA}}
\renewcommand{\headrulewidth}{0.5pt}
\renewcommand{\footrulewidth}{0.5pt}

% Section formatting
\titleformat{\section}
  {\normalfont\sffamily\Large\bfseries\color{blue!80!black}}
  {\thesection}{0.8em}{}
\titleformat{\subsection}
  {\normalfont\sffamily\large\bfseries\color{blue!70!black}}
  {\thesubsection}{0.6em}{}
\titleformat{\subsubsection}
  {\normalfont\sffamily\normalsize\bfseries\color{black}}
  {\thesubsubsection}{0.5em}{}

\titlespacing*{\section}{0pt}{10pt}{5pt}
\titlespacing*{\subsection}{0pt}{8pt}{4pt}
\titlespacing*{\subsubsection}{0pt}{5pt}{2pt}
\setlength{\parskip}{3pt}
\setlength{\parindent}{0pt}

% Code listing styling
\definecolor{codegray}{rgb}{0.96,0.96,0.96}
\definecolor{codeframe}{rgb}{0.85,0.85,0.85}
\lstset{
  backgroundcolor=\color{codegray},
  basicstyle=\ttfamily\footnotesize,
  breaklines=true,
  frame=single,
  rulecolor=\color{codeframe},
  keepspaces=true
}

\title{\sffamily\bfseries\huge HashiCorp Vault Guide\\[0.3em]
\Large Architecture, Operations, and Python (\texttt{hvac}) Microservice Integration}
\author{\textsf{The Guardian Security \& Infrastructure Team}}
\date{\textsf{Version 1.0 \textbar\ November 2026}}

\begin{document}

\maketitle
\thispagestyle{fancy}

\begin{abstract}
\noindent This guide establishes the operational and architectural standards for HashiCorp Vault deployment within the 3-layer urban monitoring platform. It details Raft integrated storage topology, AES-256-GCM encryption at rest, Shamir secret unsealing, AppRole dynamic authentication, KV-v2 versioning, Transit cryptographic acceleration, and automated Python token renewal loops via \texttt{hvac}.
\end{abstract}

\vspace{0.3cm}

\section{Architecture and Operations Overview}

\subsection{Cluster Topology}
Vault executes as a high-availability cluster utilizing \textbf{integrated storage (Raft consensus)}, eliminating external database dependencies such as Consul or ZooKeeper.

\begin{table}[htbp]
\centering
\small
\begin{tabularx}{\textwidth}{p{3.2cm}X}
\toprule
\textbf{Concept} & \textbf{Operational Mechanism} \\
\midrule
\textbf{Raft Consensus} & Nodes elect a single active leader. Writes route exclusively through the leader and replicate to followers, committing once a quorum ($n/2 + 1$ nodes) acknowledges storage. \\
\textbf{Fault Tolerance} & A 3-node cluster tolerates 1 node outage; a 5-node cluster tolerates 2 outages. Clusters require an odd node count. \\
\textbf{Active / Standby} & The elected leader serves live requests. Standby nodes forward incoming client requests to the active leader via internal mTLS and execute failover election upon heartbeat timeout. \\
\textbf{Network Ports} & Port 8200 serves the external client REST API and CLI; Port 8201 serves cluster-internal Raft replication and request forwarding. \\
\bottomrule
\end{tabularx}
\end{table}

\subsection{Security Model}
\begin{itemize}[leftmargin=*,itemsep=2pt]
  \item \textbf{TLS Transport Enforcement:} The API listener enforces HTTPS. Client microservices must pass a verified CA bundle (\texttt{VAULT\_CACERT}) to validate server hostnames. Disabling certificate verification (\texttt{verify=False}) is strictly forbidden.
  \item \textbf{Envelope Encryption at Rest:} Storage records are encrypted with AES-256-GCM using Vault's internal encryption key, which is itself wrapped by the cluster root key. Stolen disks or snapshot volumes yield only ciphertext.
  \item \textbf{Shamir Secret Unsealing:} Following restart, Vault enters a sealed state. Unsealing requires a threshold quorum of key shares (e.g., 3-of-5 threshold) generated during cluster initialization. Unseal shares must be partitioned across administrative officers and stored in HSMs or encrypted keychains.
  \item \textbf{Engines Employed:}
    \begin{itemize}
      \item \textbf{KV v2 (\texttt{secret/}):} Versioned secret storage for database credentials and service parameters.
      \item \textbf{Transit (\texttt{transit/}):} Cryptography-as-a-Service providing RSA-4096 and Ed25519 offloading without exposing private key material to client applications.
    \end{itemize}
  \item \textbf{Audit Logging:} File-based audit devices log all cryptographic operations and access attempts to \texttt{/vault/logs/audit.log}.
\end{itemize}

\subsection{AppRole Authentication Workflow}
Client microservices authenticate via AppRole:
\begin{enumerate}[leftmargin=*,itemsep=2pt]
  \item \textbf{Policy Definition:} Administrators define path-scoped HCL policies enforcing default-deny semantics.
  \item \textbf{Credential Provisioning:} The role yields a static \texttt{role\_id} (application identifier) and a dynamically injected, short-lived \texttt{secret\_id} (mounted via Docker secret).
  \item \textbf{Token Issuance:} The application exchanges credentials at \texttt{auth/approle/login}, receiving a 15-minute scoped client token.
\end{enumerate}

\section{Python (\texttt{hvac}) Integration Guide}

\subsection{Client Initialization and AppRole Login}
\begin{lstlisting}
import os, hvac

def make_client() -> hvac.Client:
    return hvac.Client(
        url=os.environ["VAULT_ADDR"],
        verify=os.environ.get("VAULT_CACERT", True),
        timeout=10,
    )

def login(client: hvac.Client) -> int:
    response = client.auth.approle.login(
        role_id=os.environ["VAULT_ROLE_ID"],
        secret_id=os.environ["VAULT_SECRET_ID"],
    )
    if not client.is_authenticated():
        raise RuntimeError("AppRole login failed: token invalid")
    return response["auth"]["lease_duration"]
\end{lstlisting}

\subsection{Background Token Renewal Daemon}
To prevent service outages upon token expiration, microservices utilize a background renewal thread:
\begin{lstlisting}
import threading, logging

class VaultSession:
    def __init__(self):
        self.client = make_client()
        self._ttl = login(self.client)
        self._stop = threading.Event()
        self._thread = threading.Thread(target=self._renew_loop, daemon=True)
        self._thread.start()

    def _renew_loop(self):
        while not self._stop.wait(max(self._ttl * 2 / 3, 1)):
            try:
                lease = self.client.auth.token.renew_self()["auth"]["lease_duration"]
                self._ttl = lease
            except Exception:
                self._ttl = login(self.client)

    def close(self):
        self._stop.set()
        self._thread.join(timeout=5)
        try:
            self.client.auth.token.revoke_self()
        finally:
            self.client.adapter.close()
\end{lstlisting}

\section{Transit Decryption Service}
Departmental RSA-4096 private keys remain sealed inside Vault. Services forward encrypted payloads in Transit format (\texttt{vault:v1:<base64>}) and receive plaintext bytes:
\begin{lstlisting}
import base64

def decrypt_for_department(client: hvac.Client, department: str, ciphertext: str) -> bytes:
    response = client.secrets.transit.decrypt_data(
        name=f"dept-{department}",
        ciphertext=ciphertext,
        mount_point="transit",
    )
    return base64.b64decode(response["data"]["plaintext"])
\end{lstlisting}

\section{Security Verification Checklist}
\begin{itemize}[leftmargin=*,itemsep=2pt]
  \item[\checkmark] TLS 1.3 enforced on listener; client CA bundle verification enabled.
  \item[\checkmark] Initial root token revoked; Shamir shares distributed across authorized officers.
  \item[\checkmark] AppRole \texttt{secret\_id} injected via runtime secrets, absent from Git and container images.
  \item[\checkmark] Path-level least-privilege policies active for \texttt{guardian-app}, \texttt{twin-app}, and \texttt{brain-app}.
  \item[\checkmark] Audit device writing structured audit records to persistent volume.
\end{itemize}

\end{document}

